\documentclass[11pt]{article}
\usepackage[a4paper,margin=25mm]{geometry}
\usepackage{amsmath,amssymb,amsthm,mathtools}
\usepackage{microtype}
\usepackage[colorlinks=true,linkcolor=blue,urlcolor=blue,citecolor=blue]{hyperref}
\newtheorem{theorem}{Theorem}
\newtheorem{lemma}{Lemma}
\newtheorem{corollary}{Corollary}
\newcommand{\PP}{\mathsf P}
\newcommand{\QQ}{\mathsf Q}
\newcommand{\EE}{\mathbb E}
\newcommand{\TV}{\mathrm{TV}}
\newcommand{\Mis}{\mathrm{Mis}}
\title{Frequentist Approximate Sufficiency of Spectral Information\\for SBM Partitions}
\author{Spectral-LRT research note}
\date{12 October 2026}
\begin{document}
\maketitle
\begin{abstract}
The partition $C$ is a fixed parameter throughout this note. No prior,
random partition, or posterior argument is used. Uniform exact recovery
from a statistic implies approximate sufficiency for bounded finite-action decisions
whose loss depends only on the partition, even when the connection matrix
is unknown. Combining this reduction with a published rowwise eigenspace
perturbation theorem proves the result for every fixed number of blocks,
uniform spectral separation, and $np\ge M\log n$ with a sufficiently large
constant $M$. A separate fixed-known-$P$ result gives ordinary Le Cam
asymptotic sufficiency. These statements do not assert equivalence of the
full experiment with unknown parameter $(C,P)$.
\end{abstract}

\section{Fixed-parameter model and the target notion}
Let $C$ be a partition of $[n]$ into $K$ nonempty blocks. Label permutations
are identified. A canonical ordering, for example by the smallest vertex in
each block, supplies indices for the probability matrix $P$. Its indices
must always be changed together with any alternative block ordering.
For each fixed $\theta=(C,P)$, independently for $i<j$,
\[
 A_{ij}\sim\operatorname{Ber}(P_{c_i c_j}),\qquad
 A_{ji}=A_{ij},\qquad A_{ii}=0.
\]
The symmetric matrix $P$ contains the actual probabilities. Set
$p=\max_{a,b}P_{ab}$. Fix $K\ge2$ and positive constants
$\pi_0,\kappa,\eta$. The parameter class satisfies
\begin{equation}\label{eq:model}
 n_a\ge\pi_0 n,\qquad
 \kappa p\le P_{ab}\le p\le1-\eta,\qquad
 \sigma_{\min}(P)\ge\kappa p.
\end{equation}
There is no convergence assumption on $P/p$ or $n_a/n$. Retain the $K$
adjacency eigenpairs with largest absolute eigenvalues, including their signs:
\[
 T(A)=(U,\Lambda),\qquad U^\top U=I_K.
\]
Use measurable deterministic tie conventions. Write
$\PP_{C,P}$ for the graph distribution and $\QQ_{C,P}$ for the induced
distribution of $T$. Neither is a distribution placed on the parameter.

For a finite action space $\mathcal D$ and any loss
$0\le L(C,d)\le1$, let
$R_{C,P}(\delta)=\EE_{C,P}L(C,\delta)$.
We say that $T$ is \emph{$\epsilon$-sufficient for partition decisions}
if, for every such loss and every possibly randomized full-data rule
$\delta_A$, there is a rule $\delta_T$ using only $T$ such that
\begin{equation}\label{eq:decision}
 \sup_{(C,P)\in\Theta_n}
 \{R_{C,P}(\delta_T)-R_{C,P}(\delta_A)\}\le\epsilon.
\end{equation}
All losses here are invariant under block-label permutations and evaluate
the partition only. This is the ordinary decision-theoretic comparison
restricted to the parameter of interest. The rule $\delta_T$ cannot use
the unknown $P$.

\section{A finite-sample reduction requiring no prior}
\begin{theorem}[Recovery implies partition-decision sufficiency]\label{thm:reduction}
Suppose a statistic-based estimator satisfies
\[
 \sup_{(C,P)\in\Theta_n}
 \PP_{C,P}\{\widehat C(T)\ne C\}\le\epsilon_n.
\]
Then \eqref{eq:decision} holds with $\epsilon=\epsilon_n$.
\end{theorem}
\begin{proof}
For the given loss, choose a deterministic minimizer
$d^*(c)\in\arg\min_{d\in\mathcal D}L(c,d)$, and use
$\delta_T(t)=d^*(\widehat C(t))$. For each fixed $(C,P)$,
\begin{align*}
 R_{C,P}(\delta_T)
 &\le \min_d L(C,d)+\PP_{C,P}(\widehat C\ne C)\\
 &\le \min_d L(C,d)+\epsilon_n\\
 &\le R_{C,P}(\delta_A)+\epsilon_n.
\end{align*}
On the correct-recovery event the action is optimal for the fixed true
partition; on its complement the additional loss is at most one.
Every full-data action has loss at least the oracle minimum. Take the
supremum over the fixed parameters. There is no averaging over $C$.
\end{proof}

\section{A general fixed-\texorpdfstring{$K$}{K} spectral premise}
\begin{theorem}[Uniform spectral exact recovery]\label{thm:recovery}
For every $q>0$ there are $M_q,n_0$, depending only on
$K,\pi_0,\kappa,\eta,q$, such that under \eqref{eq:model},
\[
 n\ge n_0,\quad np\ge M_q\log n
 \quad\Longrightarrow\quad
 \sup_{(C,P)\in\Theta_n}
 \PP_{C,P}\{\widehat C(T)\ne C\}\le n^{-q}.
\]
The estimator described below needs neither $P$ nor the population centers.
\end{theorem}

\paragraph{Population geometry.}
Let $Z$ be the membership indicator, $N=\operatorname{diag}(n_a)$,
$F=ZN^{-1/2}$, and $D_{ii}=P_{c_i c_i}$. The actual loop-free mean is
\[
 M=\EE_{C,P}A=ZPZ^\top-D.
\]
The matrix $D$ is constant within each block. Consequently
$\operatorname{col}(F)$ is an invariant subspace of $M$, and
\[
 F^\top MF=N^{1/2}PN^{1/2}-\operatorname{diag}(P_{aa}).
\]
Its smallest singular value is at least $\pi_0\kappa np-p$.
On the orthogonal complement the eigenvalues are $-P_{aa}$, each with
multiplicity $n_a-1$. Thus, eventually, the $K$ signal eigenvalues have
absolute magnitude and separation from the remaining spectrum at least
$c np$, with an upper bound $Cnp$. Both signs are allowed.

If $V$ diagonalizes the displayed $K$-dimensional restriction, the signal
eigenvectors are $U_*=FV$. Their community centers obey
\begin{equation}\label{eq:centers}
 m_a=V_{a,:}/\sqrt{n_a},\qquad
 \|m_a-m_b\|_2=\sqrt{n_a^{-1}+n_b^{-1}}\ge\sqrt{2/n},
 \qquad \|U_*\|_{2\to\infty}\le(\pi_0 n)^{-1/2}.
\end{equation}
This uses the full rank-$K$ signal subspace. It would not follow by
retaining an arbitrary smaller number of directions.

\paragraph{Published stochastic input and its specialization.}
Use Lei's generic Bernoulli eigenspace perturbation theorem
\cite[Theorem 3.4]{lei}. Apply it separately to the positive and negative
signal tails, which are consecutive eigenvalue intervals, and combine
their orthogonal alignments. Theorem 5.2 of that paper is a direct
fixed-shape SBM recovery result; we instead specialize the generic bound
to obtain uniform constants on \eqref{eq:model}.

The quantities in the generic theorem satisfy, uniformly,
\[
 \lambda_{\min}^{\rm abs},\Delta\ge cnp,\quad
 p^*,\bar p^*\le p,\quad
 \|M\|_{2\to\infty}\le p\sqrt n,\quad
 \|U_*\|_{2\to\infty}\le (\pi_0 n)^{-1/2},\quad r\le K.
\]
Choose its confidence parameter $\delta=n^{-(q+2)}$ and auxiliary exponent
$1/2$. Then $R=O(\log n)$ and
$g=O(\sqrt{np}+\log n/\log\log n)$. Substitution in that theorem gives
\begin{equation}\label{eq:rowwise}
 \inf_{W\in O(K)}\|UW-U_*\|_{2\to\infty}
 \le \frac{B_q}{\sqrt n}\sqrt{\frac{\log n}{np}},
\end{equation}
outside an event of probability at most $n^{-q}$, after increasing $n_0$.
The intermediate terms are bounded by
\[
\frac{B_q}{\sqrt n}\left\{\frac{1}{\sqrt{np}}+
\frac{\log n}{np}+\sqrt{\frac{\log n}{np}}\right\}.
\]
The additional factor
\[
 \frac{\sqrt{np}\,R^{1/4}}{np\log R}
\]
is bounded and vanishes. For the final minimum in the generic theorem,
use its mean-row-norm branch. This contributes at most a constant times
$n^{-1/2}(np)^{-1/2}$.
For $np\ge M_q\log n$ these terms yield \eqref{eq:rowwise}.
The theorem's gap requirement holds eventually. The accompanying
operator-norm control identifies the retained empirical eigenpairs with
the signal tails. No extra assumption from its Corollary 3.6
is needed. Constants depend only on the fixed bounds in \eqref{eq:model}.

\paragraph{Explicit decoder and completion of the proof.}
Construct an auxiliary graph on the vertices by connecting $i,j$ if
\[
 \|U_{i,:}-U_{j,:}\|_2\le\frac1{2\sqrt n}.
\]
Return its components if there are exactly $K$ and each has size at least
$\pi_0n$. Otherwise return a fixed admissible fallback partition. The
balance bound is a specified model-class constant, not unknown $P$.
This is a measurable function of $U$; pairwise distances are unchanged
by orthogonal alignment.

Increase $M_q$ so that the right side of \eqref{eq:rowwise} is at most
$1/(8\sqrt n)$. On that event, within-community distances are at most
$1/(4\sqrt n)$, while between-community distances are at least
$(\sqrt2-1/4)/\sqrt n>1/(2\sqrt n)$ by \eqref{eq:centers}.
Every true block is therefore a clique and there are no cross-block edges.
The components are exactly $C$. This proves Theorem~\ref{thm:recovery}.

\begin{corollary}[Unknown-$P$ partition sufficiency]\label{cor:main}
Under Theorem~\ref{thm:recovery}, the spectral statistic is
$n^{-q}$-sufficient for all bounded finite-action partition decisions. In particular
\eqref{eq:decision} tends to zero uniformly with $P$ unknown.
\end{corollary}
\begin{proof}
Apply Theorem~\ref{thm:reduction} with $\epsilon_n=n^{-q}$.
\end{proof}

\section{Ordinary asymptotic sufficiency with fixed known \texorpdfstring{$P$}{P}}
This section treats $C$ as the sole experiment parameter. The fixed matrix
$P$ is known to the statistician. Block indices must be identifiable under
the convention used to generate the model: use the canonical partition
order above, or assume correctly named-label recovery. Unlabeled recovery
alone cannot justify reconstructing a model with misassigned named rows of $P$.

\begin{theorem}[Known-$P$ experiment equivalence]\label{thm:lecam}
For a fixed-known-$P$ grouping experiment, suppose
$\sup_C\PP_{C,P}(\widehat C(T)\ne C)\le\epsilon_n$.
Then the Le Cam distance between observing $A$ and observing $T$ is at
most $\epsilon_n$.
\end{theorem}
\begin{proof}
Define the Markov kernel
\[
 \mathcal K_P(t,\cdot)=\PP_{\widehat C(t),P}(\cdot).
\]
It independently generates a new graph from the estimated partition and
the known matrix $P$. The kernel does not depend on the unknown $C$.
For each fixed $C$, with $\epsilon_C=\PP_{C,P}(\widehat C\ne C)$,
\[
 \mathcal K_P\QQ_{C,P}
 =(1-\epsilon_C)\PP_{C,P}+\epsilon_C\mathsf H_C
\]
for some probability measure $\mathsf H_C$ (arbitrary when $\epsilon_C=0$).
Hence
\[
 \sup_C\|\mathcal K_P\QQ_{C,P}-\PP_{C,P}\|_{\TV}\le\epsilon_n.
\]
The other direction has deficiency zero because $T$ is a deterministic
function of $A$. Taking the maximum of the two deficiencies proves the
claim. This is the standard experiment-comparison notion of asymptotic
sufficiency \cite{han}. No randomization of the parameter has occurred.
\end{proof}

\section{Connection to the optimal error exponent}
The following is an implication, not a new verification of an existing
end-to-end decoding theorem. If a spectral-only estimator is independently
proved to obey, uniformly over a parameter class,
\[
 \EE_{C,P}\Mis(\widehat C,C)
 \le\exp\{-(1-o(1))I_n\},
\]
then, since any wrong partition has misclassification at least $1/n$,
\[
 \PP_{C,P}(\widehat C\ne C)
 \le n\EE_{C,P}\Mis
 \le n\exp\{-(1-o(1))I_n\}.
\]
If $I_n\ge(1+\tau)\log n$ for fixed $\tau>0$, this tends uniformly to zero
and gives the same sufficiency consequences. The existing manuscript's
entire optimal-exponent argument is not independently re-certified here.

\section{Precisely what is proved}
The proved result covers every fixed $K$, balanced positive block sizes,
uniform spectral separation, and logarithmic or denser degree with a
sufficiently large coefficient. It requires no prior on $C$, no known $P$
for partition decisions, no convergence of normalized model matrices,
and no restriction to assortative or two-block graphs.

The result does not provide a sharp recovery constant, or sufficiency below
the exact-recovery regime. Exact recovery is sufficient for our reduction;
we do not assert that it is necessary for sufficiency itself. With unknown
$P$ the full $(C,P)$ experiment is not declared equivalent: the kernel in
Theorem~\ref{thm:lecam} would then use an unknown parameter and be inadmissible.
Finally the theorem concerns the information in $T$. It does not identify
a fractional Bernoulli working likelihood with the true spectral density,
prove convergence of a specific likelihood refinement, or establish MCMC
mixing.

\begin{thebibliography}{9}
\bibitem{lei} Lihua Lei (2020).
\emph{Unified $\ell_{2\to\infty}$ Eigenspace Perturbation Theory for Symmetric
Random Matrices}. arXiv:1909.04798v2, Theorems 3.4 and 5.2.
\url{https://arxiv.org/pdf/1909.04798}.
\bibitem{han} Yanjun Han (2021).
\emph{Statistical decision theory: model distance and equivalence},
Lecture 4, pp.\ 11--18.
\url{https://web.stanford.edu/class/ee378c/lecture4_annotated.pdf}.
\end{thebibliography}
\end{document}
